\chapter{Weitere Anwendung und Verbreitung}

\section{Go}

Bei Go wird die Spielstärke traditionell in Kyū-Graden für Schüler und Dan-Graden für Meister angegeben. Die Ermittlung dieser Spielstärke basiert innerhalb der \textit{European Go Federation} und bei vielen Go-Servern im Internet auf einem von Elo abgeleiteten System, welches Kyū- und Dan-Grade wie folgt abbildet:

\begin{table}[htbp]
    \centering
    \begin{tabular}{l | c | l}
        \toprule
        \textbf{kyū/dan} & \textbf{Elo} & \textbf{Spielstärke und -erfahrung} \\
        \midrule
        weltbeste & $5185$ & AlphaGo Zero auf einem TPU-v2-Modul mit \\
        9p-KI & & 180 TFLOPS \cite{silver2017} \\
        weltbester & $3830$ & Shin Jin-Seo, weltbester Go-Spieler \\
        9p-Spieler & & \\
        1p - 9p & ab circa $2600$ & professioneller Go-Spieler (aus Japan, Korea oder \\
        & & China), der stärker als ein Amateur-6dan spielt \\
        4d - 7d & ab $2350$ & einer der besten Spieler seines Landes \\
        1d - 3d & $2050 - 2349$ & sehr guter Club-Spieler \\
        4k - 1k & $1650 - 2049$ & guter Club-Spieler \\
        9k - 5k & $1150 - 1649$ & regelmäßiger Club-Spieler \\
        13k - 10k & $750 - 1149$ & Club-Spieler \\
        17k - 14k & $350 - 749$ & regelmäßiger Hobby-Spieler \\
        21k - 18k & $0 - 349$ & Hobby-Spieler \\
        24k - 22k &  & einige Partien gegen Nicht-Anfänger gewonnen \\
        27k - 25k &  & einige Partien gegen Anfänger gewonnen \\
        29k - 28k &  & einige Partien gespielt \\
        30k &  & Regeln verstanden, aber noch keine Partie gespielt \\
        \bottomrule
    \end{tabular}
    \caption{Die European Go Federation leitet die Spielstärke aus der Elo ab.}
    \label{tab:go_elo}
\end{table}




\section{Fußball}

Die FIFA-Weltrangliste der Frauen wird seit 2003 offiziell mit einem adaptierten Elo-System ermittelt. Seit 2018 wird die FIFA-Weltrangliste für Männer auch auf ein adaptiertes Elo-System umgestellt. \cite{fifa2018weltrangliste}

Eine länger bestehende inoffizielle Adaption des Elo-Systems für Männernationalmannschaften im Fußball sind die World Football Elo Ratings. Inoffizielle Elo-Ratings werden auch für Fußball-Clubs vorgenommen.


\section{Tischtennis}

\subsection{In der Schweiz}

Swiss Table Tennis nutzt seit der Saison 2010/2011 eine etwas modifizierte Elo-Formel zur Berechnung von Wertungspunkten
\[
    E_{A} = \frac{1}{1 + 10^{(R_{B} - R_{A}) / 200}}
\]
\begin{itemize}
    \item[] $E_{A}$: Erwarteter Punktestand für Spieler $A$.
    \item[] $R_{A}$: bisherige Punkte-Zahl von Spieler $A$.
    \item[] $R_{B}$: bisherige Punkte-Zahl von Spieler $B$.
\end{itemize}
Der Erwartungswert von $A$ beträgt nun $E_{A} \cdot 100 \%$. Die neue Punkte-Zahl von Spieler $A$ ist
\[
    R'_{A} = R_{A} + 15 \cdot (S_{A} - S_{E})
\]
\begin{itemize}
    \item[] $S_{A}$: tatsächlich gespielter Punktestand ($1$ für jeden Sieg, $0$ für jede Niederlage, Remis ist im Tischtennis nicht möglich).
\end{itemize}


\subsection{In Deutschland: der TTR-Wert}

In Deutschland wird nach einem analogen System für jeden Aktiven ein \textbf{TTR-Wert} errechnet \cite{tischtennis2020heft11}. Hier wird die Wertungsdifferenz durch $150$ geteilt.
\[
    E_{A} = \frac{1}{1 + 10^{(R_{B} - R_{A}) / 150}}
\]


\section{Scrabble}

Für weltweites Scrabble (Global Scrabble) wird eine Elo-Rangliste von der World English-language Scrabble Players' Association (WESPA) geführt. Auf Rang 1 dieser Elo-Rang-liste liegt der Neuseeländer Nigel Richards ($2156$ Elo-Punkte, Stand 17. Oktober 2020).

Seit 2009 wird auch für das deutschsprachige Scrabble eine Elo-Rangliste geführt -- basierend auf Turnieren ab dem Jahr 2005. Im Juni 2022 führte der Deutsche Timon Boerner mit $1718$ Elo-Punkten die Liste an, in der $213$ Personen aufgeführt sind.


\section{League of Legends}

Im MOBA League of Legends, einer Liga eines Computer-Strategiespiels, wurde ebenfalls das Elo-System bei gewerteten Spielen verwendet. Inzwischen wurde es durch das Ligasystem ersetzt, dem aber noch das Elo-System zu Grunde liegt.

Bei einem Sieg bekommt man \textit{League Points} (LP), bei einer Niederlage werden LP abgezogen. Bis Ende 2020 musste man bei $100$ LP ein Best of three bzw. five gewinnen, um eine Liga aufzusteigen.


\section{Tennis-Elo-Wertungssystem}

Das Tennis-Elo-Wertungssystem berechnet das relative Leistungsniveau von Tennisspielern im Herreneinzel. Nach jedem Spiel erhält der Gewinner Punkte vom Verlierer. Die Anzahl der Punkte hängt von den relativen Wertungen der Spieler ab -- ein Sieg über einen Spieler mit höherer Wertung führt zu einer größeren Wertungssteigerung als ein Sieg über einen Spieler mit niedrigerer Wertung. Die Wertungen liegen bei Profispielern typischerweise zwischen $1500$ und über $2400$. Ein Unterschied von $100$ Punkten bedeutet, dass der Spieler mit höherer Wertung eine Gewinnchance von etwa $64 \%$ hat. $200$ Punkte bedeuten $76 \%$, $300$ Punkte bedeuten $85 \%$, $400$ Punkte bedeuten $91 \%$ und $500$ Punkte bedeuten $95 \%$ Gewinnwahrscheinlichkeit. Die Wertungen werden nach jedem Spiel dynamisch angepasst und spiegeln die aktuelle Leistung wider. Elo-Wertungen werden ausschließlich anhand von Spieldaten von Turnieren der Association of Tennis Professionals (ATP) der Herren berechnet. Sie werden insbesondere im Bereich der Tenniswetten verwendet.